\chapter{Scientific Problem}
\label{section:scientificProblem}

\section{Problem definition}
\label{section:problemDefinition}

\paragraph{Informal description.} Smartwatches and chest patches can record a single-lead ECG during everyday life, but these recordings are hard to use for three reasons:
\begin{itemize}
	\item \textbf{Noise.} The person is moving, and wearable sensors have many sources of error, so the signal is often disturbed \cite{bayoumy2021smart}.
	\item \textbf{Volume.} A wearable records all day, every day, which is far more data than a clinician can read.
	\item \textbf{Every patient is different.} ECG shapes vary from person to person, so models trained on other people make more mistakes on a new patient \cite{dechazal2004automatic,kiranyaz2016realtime}.
\end{itemize}

Our idea is a \emph{cardiac digital twin}: a model of one patient's heart that is kept up to date with that patient's own data \cite{corralacero2020digital}. Existing cardiac twins are built from hospital imaging \cite{gillette2021framework}; ours uses only wearable data. It is a simple mathematical model of a heartbeat \cite{mcsharry2003dynamical}, adjusted to the patient's own normal beats. When a new beat looks very different from what the twin expects, it is flagged as unusual for this patient. As in the course brief, five agents share the twin: Signal Monitoring, Arrhythmia Detection, Physiological Prediction, Twin Calibration and Clinical Recommendation. The result is a short report for a doctor; the system does not make diagnoses.

\paragraph{Why an intelligent algorithm?} Fixed rules fail when the signal is noisy or the patient is atypical, while deep networks trained on large datasets classify heart rhythms from single-lead ECG as well as cardiologists \cite{hannun2019cardiologist}. Learning is also needed to clean the signal, forecast heart rate and fit the twin to each patient.

\section{Formal problem statement}
\label{section:formal}

The input is a noisy ECG signal from a wearable, possibly together with heart rate and motion:
\begin{equation}
\tilde x(t) = x(t) + \eta(t),
\end{equation}
where $x$ is the true heart signal and $\eta$ is the noise. From $\tilde x$ we cut out the individual heartbeats $b_1, b_2, \dots$

The twin describes a heartbeat by its five waves (P, Q, R, S, T), each with a position, a height and a width, i.e.\ 15 numbers $\vartheta_p$ for patient $p$ \cite{mcsharry2003dynamical}. Calibrating the twin means choosing $\vartheta_p$ so that the beats it generates match the patient's normal beats. For each new beat $b_k$ the twin predicts a beat $\hat b_k$, and the \emph{deviation score}
\begin{equation}
d_k = \| b_k - \hat b_k \|
\end{equation}
measures how unusual the beat is for this patient. A typical beat gives a small score, while a beat the twin cannot reproduce, for example an extra beat coming from the ventricles, gives a large one. Table~\ref{tab:subproblems} lists the input and output of each agent.

The agents work together as a pipeline (Figure~\ref{fig:graphicalAbstract}). The Signal Monitoring agent cleans the signal and marks segments that are too noisy to use. The Arrhythmia Detection agent labels every beat, using the deviation score as an extra hint. The Physiological Prediction agent forecasts the heart rate. The Twin Calibration agent keeps the twin up to date, using only clean beats that were labelled normal. Finally, the Clinical Recommendation agent collects these results into a short report and highlights anything a doctor should look at.

\begin{table}[htbp]
\caption{Input and output of each agent.}
\label{tab:subproblems}
\centering
\small
\setstretch{1.1}
\begin{tabularx}{\textwidth}{@{}>{\raggedright\arraybackslash}p{4.0cm}>{\raggedright\arraybackslash}p{4.4cm}>{\raggedright\arraybackslash}X@{}}
\toprule
\textbf{Agent} & \textbf{Input} & \textbf{Output} \\
\midrule
Signal Monitoring & raw ECG & cleaned ECG and a signal-quality score \\
Arrhythmia Detection & heartbeats and their deviation scores & a label for every beat: normal (N), supraventricular (S), ventricular (V), fusion (F) or unknown (Q) \cite{dechazal2004automatic} \\
Physiological Prediction & recent heart rate & heart-rate forecast \\
Twin Calibration & clean normal beats & updated twin $\vartheta_p$ and deviation scores $d_k$ \\
Clinical Recommendation & outputs of all agents & report for the clinician \\
\bottomrule
\end{tabularx}
\end{table}

\section{Graphical abstract}
\label{section:graphicalAbstract}

\begin{center}
\setstretch{1.0}
\resizebox{\textwidth}{!}{%
\begin{tikzpicture}[
  font=\footnotesize,
  box/.style={draw, rounded corners=3pt, align=center, inner sep=4pt, minimum height=12mm},
  io/.style={box, fill=gray!12, draw=gray!60!black, text width=20mm},
  agent/.style={box, fill=blue!7, draw=blue!55!black, text width=31mm},
  twinbox/.style={box, fill=green!10, draw=green!45!black, text width=50mm, minimum height=15mm},
  human/.style={box, fill=orange!12, draw=orange!70!black, text width=18mm},
  arr/.style={-{Stealth[length=2.2mm]}, semithick},
  lab/.style={font=\scriptsize, inner sep=1.5pt},
]
\node[io] (wear) at (0,0) {\textbf{Wearable}\\ECG, heart rate,\\motion};
\node[agent] (mon) at (3.5,0) {\textbf{Signal Monitoring}\\noise removal,\\quality check};
\node[agent] (pred) at (7.7,1.25) {\textbf{Physiological}\\\textbf{Prediction}\\heart-rate forecast};
\node[agent] (arr) at (7.7,-1.25) {\textbf{Arrhythmia}\\\textbf{Detection}\\label for every beat};
\node[agent] (rec) at (11.9,0) {\textbf{Clinical}\\\textbf{Recommendation}\\report};
\node[human] (clin) at (15.8,0) {\textbf{Clinician}\\human in\\the loop};
\node[agent] (cal) at (3.5,-3.7) {\textbf{Twin Calibration}\\fits the twin to\\normal beats};
\node[twinbox] (twin) at (9.5,-3.7) {\textbf{Cardiac digital twin}\\model of the patient's heartbeat,\\rhythm and history};

\draw[arr] (wear) -- node[lab, above]{ECG} (mon);
\draw[arr] (mon.east) -- node[lab, above left, pos=0.5]{clean signal} (pred.west);
\draw[arr] (mon.east) -- (arr.west);
\draw[arr] (mon.south) -- node[lab, right]{clean beats} (cal.north);
\draw[arr] (cal) -- node[lab, above]{update} (twin);
\draw[arr] (arr.south |- twin.north) -- node[lab, left]{deviation score} (arr.south);
\draw[arr] (pred.east) -- node[lab, above right, pos=0.35]{forecast} (rec.west);
\draw[arr] (arr.east) -- node[lab, below right, pos=0.35]{beat labels} (rec.west);
\draw[arr] (rec.south |- twin.north) -- node[lab, right]{twin state} (rec.south);
\draw[arr] (rec) -- node[lab, above]{report} (clin);
\draw[arr, dashed] (clin.south) |- node[lab, right, pos=0.3]{feedback} (twin.east);
\end{tikzpicture}}
\captionof{figure}{Graphical abstract: five agents share a model of the patient's heart (the twin); the clinician's feedback updates it.}
\label{fig:graphicalAbstract}
\end{center}

